\documentclass[12pt]{article}
\usepackage{fontspec}
\setmainfont[AutoFakeBold=2,AutoFakeSlant=.2]{SimSun}
\usepackage[a4paper,margin=2.5cm]{geometry}
\usepackage{amsmath,array,longtable,hyperref}
\hypersetup{hidelinks,pdfauthor={}}
\XeTeXlinebreaklocale "zh"\XeTeXlinebreakskip=0pt plus 1pt
\setlength{\parindent}{2em}\setlength{\parskip}{0.35em}
\setlength{\tabcolsep}{4pt}\renewcommand{\arraystretch}{1.2}
\renewcommand{\tablename}{表}\renewcommand{\refname}{参考文献}
\renewcommand{\abstractname}{摘要}
\newcommand{\toprule}{\noalign{\hrule height .8pt}}
\newcommand{\midrule}{\noalign{\hrule height .4pt}}
\newcommand{\bottomrule}{\noalign{\hrule height .8pt}}
\Urlmuskip=0mu plus 1mu\emergencystretch=2em
\makeatletter
\def\ps@custom{\def\@oddhead{}\def\@evenhead{}\def\@oddfoot{\footnotesize 市场状态变化下的量化投资策略\hfil\thepage}\let\@evenfoot\@oddfoot}
\makeatother\pagestyle{custom}
\hypersetup{pdftitle={市场状态变化下的量化投资策略},pdfsubject={从量价选股到多资产风险预算的成本后验证}}
\begin{document}
\begin{center}{\Large\bfseries 市场状态变化下的量化投资策略}\\[.5em]{\normalsize 从量价选股到多资产风险预算的成本后验证}\\[.5em]{\small 2026年10月10日}\end{center}
\begin{abstract}
本研究从历史沪深300股票池的量价选股出发，比较财务因子、趋势、状态切换、风险控制、新闻主题和注意力模型，再扩展到股票、标普500、黄金及国债的共享现金账户。E00股票基线在2024年收益39.49\%、2025年8.05\%，表明跨年度表现分化。最终风险预算候选RISK-B1在2024—2025年标准成本下累计收益27.80\%、零现金夏普2.858、最大回撤3.38\%；同对象固定25\%配置为46.13\%、2.642及5.55\%。B1提高单位波动收益并降低回撤，代价是绝对收益下降。2026年175日现金起点补充检验中，B1标准收益2.96\%、夏普0.446、最大回撤8.76\%，固定配置为0.52\%、0.122及10.22\%。结果为已查看历史上的探索比较。本文提供实际成交、费用及公司行为约束下的可复核策略与比较过程，不宣称全局最优或实盘保证。 股票、新闻、深度学习及配置分支按信息来源与评价时期分别比较；2023为开发验证，2024—2025为已观察探索，2026为部分年度的事后补充检验。注意力对配对MLP的局部优势与新闻当前预测的负结果均保留。主策略复现与扩展研究的复现范围分别列明。
\end{abstract}
\noindent\textbf{关键词：}量化选股；多资产配置；风险预算；成本后夏普；宏观主题NLP

所有已完成与未完成路线分别标记，策略参数保持冻结。


\section{研究问题、经济机制与实验范围}
策略目标是在可执行的长仓账户中提高成本后夏普。股票选股回答相对买谁，资产配置回答把多少资本交给股票、海外股票、黄金及利率资产；现金始终可持有。前期曾要求新模块同时改善弱年收益、回撤与成本，后期明确改为夏普优先，仍同时披露收益与回撤。前后目标变化不能混写成一套从未改变的验收标准。


研究问题不是“能否堆叠更多模型”，而是：已有量价关系失效时，增加信息、改变期限、识别条件或调整风险，能否改善成本后收益且不产生不可接受的回撤与执行成本？金融机制先于算法。

量价信号依赖参与者行为在价格与成交中留下可重复关系；财务补充经营与估值信息；趋势捕捉持续性；区制尝试识别关系变化；风险配置改变资本承担；NLP尝试描述信息与预期。这些机制可能重叠，也可能适用于不同期限。

表\ref{tab:route}给出对应路线或模型的比较结果。


{\small
\begin{longtable}{@{}p{2.5cm}p{5.5cm}p{6.0cm}@{}}
\caption{实验路线及完成范围}\label{tab:route}\\
\toprule
路线 & 实际做了什么 & 证据结论 \\
\midrule\endfirsthead
\multicolumn{3}{c}{表\thetable（续）}\\\toprule
路线 & 实际做了什么 & 证据结论 \\
\midrule\endhead
量价与简单基准 & 等权、规则、E00及E10账户 & E00强年较好，弱年不足 \\
财务与公开因子 & 轻量联合、Qlib适配、JF05/JF20 & 未证明足够成本后增量 \\
趋势与切换 & T01、S01、独立50/50双袖 & 部分年份改善，未达全套门槛 \\
风险仓位 & 12\%波动目标与固定仓位对照 & 有风险作用，动态版未达门槛 \\
HMM与序列 & RG01专家、SQ01储备池与Ridge & 历史账户完成；SQ01已补2026账户 \\
宏观主题NLP & 六主题、价格/文本/延迟文本对照 & 当前指数目标未证明预测增量 \\
复杂配置与RL & BLO、DRO、HRP、强化学习等讨论 & 未实现完整检验，不算失败实验 \\
\bottomrule\end{longtable}}
BLO、DRO、HRP和宏观强化学习在前期被提出或讨论，未形成完整候选账户实验；债券已通过国债ETF进入后期多资产验证，直接期货未实施。研究范围不能把讨论记录等同模型实现；多资产可用也不代表应机械转入债券。

各轮比较记录各自冻结规则；后期固定配置曾进行132组预设候选搜索。原股票升级采用多目标门槛，后期投资目标明确改为夏普优先，目标变更与多次选择偏差须披露。失败只否决当前数据、目标与实现下的升级，不否定整个方法类别。


\section{数据来源、可得性与账户口径}
掘金数据提供历史行情、复权相关字段、月度沪深300候选池及账户所需输入。历史129个月股票池的证券并集为637只；并集并不是每天可以买的股票。信号使用复权价格，账户使用原价，不能把今日成分回填历史。

比赛邮件已交付8653条业绩预告，证券代码桥接、公告原始与修订版本链仍有缺口；它们不是E00模型输入。iFinD提供新闻检索标题、日期、URL与正文片段，2025现有2349条片段覆盖245个文章日期；这是抽样检索结果，并非完整历史新闻库。

财务与估值用于财务候选，公告时间严格早于信号日、估值取此前交易日，并记录陈旧度与缺失；历史修订版本和行业适用性仍有限。Tushare的权限验证与宏观小样本不等于完整宏观训练序列。2026真实新闻缺口不以零值填充。

股票基线账户起始100万元、股票与现金、现金收益假设为0；80目标持仓、前120保留缓冲、每5交易日检查、单股买入上限3\%。信号收盘产生，下一交易日开盘执行，考虑整手、T+1、停牌、涨跌停及公司行为。有限日线不能确定真实盘口排队和部分成交。

2026标准佣金假设0.03\%、每笔至少5元，另含过户费与卖方印花税；单边滑点5基点。原压力按冻结适配器把交易费用加倍、滑点10基点，股息税另行按持有期处理。股息应收、税负与送转独立记账，不从复权价格再次计入账户。

供应商来源、原发布机构与镜像获取渠道分别披露；计划派息不等于券商实收，学校访问权限不等于可以公开整个数据库。最终交付重点重放RISK-B1与固定25\%共享账户；扩展诊断的完整输入与复现范围在交付清单中另列。


最终配置使用两个职责不同的ETF输入：风险代理历史报价来自Tushare \texttt{fund\_daily}，原成交额按千元转换成人民币元；实际账户ETF开收盘报价来自东方财富 \texttt{push2his}第三方行情接口。二者保存来源字段及响应哈希，不默认为同一家供应商的同一版本。股票核心行情及历史成分来自掘金；权益事件结合初始供应商记录、BaoStock日期与发行人公告核实补丁。新闻、比赛邮件及资金资料分别用于对应探索或覆盖核查，未混入最终E00与风险预算的训练输入。

\section{基线模型、时间切分与候选实现}
E00为20项价格/交易特征的LightGBM，包含收益、趋势、波动、量价、流动性及beta等。每月使用36个月滚动训练窗，15叶、学习率0.05、叶内最少100样本、L2为5、100轮、固定种子20261006。训练目标为下一日开盘至第六日开盘的复权收益，转为当日截面百分位减0.5，各日期总权重相同。

月度训练只使用月初前一全局交易日已成熟的标签。2026九个月按冻结在线日程允许较早2026成熟样本进入后月训练；这不是固定只用2020—2024训练，也不是用整个2026结果一次性训练。预测股票不因未来标签尚未成熟而被删除。

E10增加五个市场/候选池状态变量。P20把目标改为t+1开盘至t+21开盘，仍每五日检查持仓；预测期限与调仓频率不同。JF20在同期限P20上增加14项财务特征。Qlib路线是公开因子与集成思路的适配，不是原论文收益的精确复现\cite{ref2}。

T01取60日、120日收益与跳过近期的中期动量的截面排名平均。S01在市场20日收益及相对60日均线均为正时使用T01，其他使用E00。双袖组合为初始各50万元独立账户，无跨袖再平衡；整手与最低费用存在，不能直接把两个百万账户收益平均。

R01动态仓位为 min(1, 0.12/预测年化波动)，另以开发期固定仓位0.5964844995作对照。RG01沿用冻结HMM状态概率混合两组专家；SQ01采用固定储备池状态与Ridge。隐状态是统计状态，不自动等于金融意义的牛熊\cite{ref3}。

训练目标及风险暴露分别表示为

\begin{equation}y_{i,t}^{(h)}=\widetilde O_{i,t+h+1}/\widetilde O_{i,t+1}-1,\quad z_{i,t}=\operatorname{RankPct}_i(y_{i,t}^{(h)})-1/2.\end{equation}
其中 O\textasciitilde{} 为复权开盘价，h=5 为E00目标、h=20 为P20目标；标签终点须不晚于训练截止日。风险仓位为

\begin{equation}a_t=\min\{1,0.12/\widehat\sigma_t\}.\end{equation}
其中 $\widehat{\sigma}_t$ 为共同目标股票篮子过去20个交易日等权收益的年化波动估计。零波动与无效输入按冻结规则处理。账户净资产按原价与现金流记账：

\begin{equation}V_t=C_t+\sum_i q_{i,t}P_{i,t}+D_t-T_t.\end{equation}
其中 C(t) 为现金、D(t) 为股息应收、T(t) 为估计股息税负；复权仅用于信号，不替代账户分红与股数变动。


\section{历史样本回测与候选取舍}
评价口径：累计收益为期末NAV/初始资金-1；年化收益按交易日数折算；最大回撤从包含初始资金的历史峰值计算。不同股票暴露的收益不能直接当作alpha。2026仅部分年度，正文不将其收益年化后当作全年实际收益。股票历史表来自既有账户日表重汇总；后期多资产候选另做共享现金账户回放，二者评价区间与池子分别注明。

先以共同历史股票池和100万元现金起点比较量价、财务及公开因子路线，再开展期限、趋势、风险和区制实验。表\ref{tab:history}汇总标准成本485日账户结果，表\ref{tab:annual}说明主要信号的年度分化。


{\small
\begin{longtable}{@{}lrrrrr@{}}
\caption{2024—2025年标准成本账户比较（收益与回撤，\%）}\label{tab:history}\\
\toprule
模型 & 2024年 & 2025年 & 累计收益 & 年化收益 & 最大回撤 \\
\midrule\endfirsthead
\multicolumn{6}{c}{表\thetable（续）}\\\toprule
模型 & 2024年 & 2025年 & 累计收益 & 年化收益 & 最大回撤 \\
\midrule\endhead
B00 & 12.95\% & 8.55\% & 22.61\% & 11.17\% & -8.64\% \\
B01 & 21.85\% & 8.58\% & 32.30\% & 15.65\% & -13.77\% \\
E00 & 39.49\% & 8.05\% & 50.72\% & 23.76\% & -8.72\% \\
E10 & 36.83\% & 6.41\% & 45.60\% & 21.56\% & -9.12\% \\
Q-FIN & 31.51\% & 5.46\% & 38.69\% & 18.52\% & -8.45\% \\
Q-LIGHT & 36.96\% & 2.40\% & 40.24\% & 19.21\% & -8.48\% \\
Q-MONTH & 31.80\% & 5.37\% & 38.88\% & 18.61\% & -8.63\% \\
JF05 & 38.72\% & 0.78\% & 39.80\% & 19.02\% & -9.33\% \\
P20 & 38.79\% & 7.50\% & 49.20\% & 23.11\% & -8.29\% \\
JF20 & 37.27\% & 4.88\% & 43.97\% & 20.85\% & -8.82\% \\
T01 & 12.43\% & 17.34\% & 31.93\% & 15.48\% & -18.64\% \\
S01 & 28.45\% & 13.48\% & 45.77\% & 21.63\% & -16.70\% \\
RG01 & 47.99\% & 5.27\% & 55.80\% & 25.91\% & -8.75\% \\
SQ01 & 21.17\% & 5.96\% & 28.39\% & 13.87\% & -9.04\% \\
R01-D & 29.41\% & 2.45\% & 32.58\% & 15.78\% & -8.71\% \\
R01-S & 20.85\% & 3.91\% & 25.57\% & 12.56\% & -5.05\% \\
\bottomrule\end{longtable}}
{\footnotesize\noindent 注：统一区间2024-01-02至2025-12-31，共485个交易日；起始资金100万元。年化收益采用252交易日约定，不采用自然日CAGR。Q-FIN为轻量财务联合，Q-LIGHT为轻量公开因子适配集成，Q-MONTH为月度完整样本版本。R01-D/S分别为动态/固定风险仓位。模型容量、更新频率与风险暴露有差异，本表不识别单个组件的因果贡献。\par}

{\small
\begin{longtable}{@{}lrrr@{}}
\caption{量价与趋势策略年度收益（标准成本，\%）}\label{tab:annual}\\
\toprule
模型 & 2023年 & 2024年 & 2025年 \\
\midrule\endfirsthead
\multicolumn{4}{c}{表\thetable（续）}\\\toprule
模型 & 2023年 & 2024年 & 2025年 \\
\midrule\endhead
E00 & 6.76\% & 39.49\% & 8.05\% \\
T01 & -2.76\% & 12.43\% & 17.34\% \\
S01 & -0.02\% & 28.45\% & 13.48\% \\
\bottomrule\end{longtable}}
{\footnotesize\noindent 注：2023年为独立现金起点的开发回溯；2024年现金重置，2025年继承2024年持仓。各年度回撤锚点不与两年累计回撤混用。\par}
同池等权B00的平均现金约43.34\%，受整手与资金规模限制，因此不是与E00同风险暴露的纯alpha对照。增加财务或公开量价因子后的不同版本同时改变容量或更新频率，不能把全部收益差归因于单一信息组。

新增模块的取舍采用多目标门槛，而非单年收益排名。JF20相对P20的标准成本两年年化差为-2.264个百分点、2025年差为-2.622个百分点，均低于原+2个百分点要求；2025年仅4/12个月胜出，低于7/12要求。RG01的两年累计收益优于E00，但2025年收益差约-2.78个百分点，未解决弱年目标。以上结论来自预先登记的门槛，不以之后发现的局部优势反改取舍。

财务指标变化频率与收益预测期限不一致、区制专家样本分割以及原有风格关系的变化，是合理解释，但上述账户结果本身不构成因果识别。2023为开发回溯，2024—2025为多次查看的探索期；年度表现相反仅提供互补研究线索，不足以证明事后组合有效。


\section{2026固定规则回测与成本敏感性}
模型编码说明：E00为量价基线；E10为量价加五项市场/候选池状态；T01为中期趋势；S01为趋势条件切换；P20为20日预测目标；R01-D/S为动态/固定风险仓位；双袖为E00与T01两个独立账户。

2026账户统一从首个可用交易日2026-01-05以100万元现金起点运行，截至2026-09-21，共175日。表\ref{tab:recent}同时列出标准/原压力收益及各自回撤；不报告“全年2026收益”。


{\small
\begin{longtable}{@{}lrrrrr@{}}
\caption{2026年固定参数账户表现及成本压力（\%）}\label{tab:recent}\\
\toprule
模型 & 标准收益 & 压力收益 & 成本拖累 & 标准回撤 & 压力回撤 \\
\midrule\endfirsthead
\multicolumn{6}{c}{表\thetable（续）}\\\toprule
模型 & 标准收益 & 压力收益 & 成本拖累 & 标准回撤 & 压力回撤 \\
\midrule\endhead
E00 & -4.31\% & -6.00\% & 1.69 & -8.56\% & -9.78\% \\
T01 & -4.49\% & -5.48\% & 0.99 & -14.06\% & -14.82\% \\
S01 & -1.86\% & -3.72\% & 1.85 & -11.08\% & -12.32\% \\
E10 & -5.48\% & -7.51\% & 2.03 & -10.73\% & -10.83\% \\
R01-D & -3.53\% & -5.29\% & 1.77 & -6.79\% & -8.20\% \\
R01-S & -2.54\% & -3.73\% & 1.19 & -4.61\% & -5.44\% \\
双袖组合 & -4.28\% & -5.98\% & 1.70 & -9.58\% & -10.62\% \\
P20 & -3.16\% & -5.31\% & 2.15 & -10.27\% & -11.47\% \\
\bottomrule\end{longtable}}
{\footnotesize\noindent 注：区间2026-01-05至2026-09-21，共175个交易日，非全年。成本拖累=标准收益-压力收益，单位为百分点；回撤为负数。压力情景沿用原冻结费用适配器，不是简单把净收益扣费。E00/T01双袖初始各50万元独立记账，无跨袖再平衡。\par}
表\ref{tab:risk}列出风险暴露和交易活动，表\ref{tab:costs}进一步区分费用、股息税及价格滑点。


{\small
\begin{longtable}{@{}lrrrrr@{}}
\caption{2026年标准成本账户风险暴露及交易特征}\label{tab:risk}\\
\toprule
模型 & 年化波动(\%) & 零现金Sharpe & 股票暴露(\%) & 交易笔数 & 盈利月份 \\
\midrule\endfirsthead
\multicolumn{6}{c}{表\thetable（续）}\\\toprule
模型 & 年化波动(\%) & 零现金Sharpe & 股票暴露(\%) & 交易笔数 & 盈利月份 \\
\midrule\endhead
E00 & 12.33\% & -0.453 & 81.40\% & 1608 & 5/9 \\
T01 & 14.31\% & -0.391 & 75.22\% & 1139 & 4/9 \\
S01 & 14.68\% & -0.112 & 76.58\% & 1747 & 5/9 \\
E10 & 14.00\% & -0.510 & 81.89\% & 1803 & 4/9 \\
R01-D & 8.19\% & -0.590 & 62.56\% & 1796 & 5/9 \\
R01-S & 6.08\% & -0.578 & 42.01\% & 1295 & 3/9 \\
双袖组合 & 9.69\% & -0.602 & 65.43\% & 2064 & 4/9 \\
P20 & 14.63\% & -0.243 & 80.46\% & 1899 & 3/9 \\
\bottomrule\end{longtable}}
{\footnotesize\noindent 注：波动和Sharpe按175个日收益以252交易日年化，仅为短样本描述，不代表未来全年水平。股票暴露为每日股票市值/NAV的算术平均；Sharpe以零现金收益为参考，并非以市场指数为基准。交易笔数计实际成交记录，双袖计两账户之和。2026-09为截至21日的部分月。\par}

{\small
\begin{longtable}{@{}lrrr@{}}
\caption{2026年标准成本费用与滑点记录（元）}\label{tab:costs}\\
\toprule
模型 & 交易费用 & 已实现股息税 & 滑点成本 \\
\midrule\endfirsthead
\multicolumn{4}{c}{表\thetable（续）}\\\toprule
模型 & 交易费用 & 已实现股息税 & 滑点成本 \\
\midrule\endhead
E00 & 11,409.20 & 1,558.01 & 6,866.65 \\
T01 & 7,552.07 & 1,351.77 & 3,985.79 \\
S01 & 12,616.73 & 1,361.75 & 7,855.72 \\
E10 & 13,005.21 & 1,380.40 & 8,078.96 \\
R01-D & 11,789.18 & 903.79 & 5,770.91 \\
R01-S & 8,166.78 & 743.87 & 3,463.10 \\
双袖组合 & 12,482.07 & 1,312.72 & 4,513.05 \\
P20 & 13,751.22 & 1,816.85 & 8,492.41 \\
\bottomrule\end{longtable}}
{\footnotesize\noindent 注：交易费用包含冻结适配器中的佣金、过户费与印花税。滑点成本已经体现在成交价中，不再从净收益重复扣除；股息税单列。该表不是不改变持仓路径的“毛收益还原”。\par}
S01相对E00的标准收益提高约2.44个百分点，但最大回撤加深约2.52个百分点；P20的收益损失较少，回撤却更深。固定仓位与动态仓位的比较还包含实际风险暴露差异，不能以损失较少直接宣布风险预测或选股alpha有效。

同期沪深300价格指数累计收益约-3.78\%；E00与其收益的描述性差额约-0.53个百分点。该指数不含股息且不承担账户交易费用，也不是同仓位可执行基准，本稿不把该差额称为风险调整alpha。

表\ref{tab:months}和表\ref{tab:months-risk}逐月列示收益路径，有助于识别改善是否集中于少数月份。


{\small
\begin{longtable}{@{}lrrrr@{}}
\caption{2026年主要策略标准成本月度收益（\%）}\label{tab:months}\\
\toprule
月份 & E00 & T01 & S01 & P20 \\
\midrule\endfirsthead
\multicolumn{5}{c}{表\thetable（续）}\\\toprule
月份 & E00 & T01 & S01 & P20 \\
\midrule\endhead
2026-01 & -0.42\% & 5.03\% & 5.03\% & -0.71\% \\
2026-02 & 0.49\% & 3.24\% & -0.29\% & 0.40\% \\
2026-03 & -2.33\% & -7.58\% & -4.81\% & -3.48\% \\
2026-04 & 2.09\% & 4.49\% & 3.98\% & 3.97\% \\
2026-05 & -2.62\% & -1.81\% & -1.50\% & -3.19\% \\
2026-06 & 0.15\% & -0.15\% & 0.63\% & -1.74\% \\
2026-07 & -3.67\% & -5.65\% & -6.43\% & -3.20\% \\
2026-08 & 2.03\% & 1.68\% & 2.03\% & 5.48\% \\
2026-09 & 0.06\% & -3.02\% & 0.06\% & -0.33\% \\
\bottomrule\end{longtable}}
{\footnotesize\noindent 注：每月收益=本月最后可用交易日NAV/上月最后可用交易日NAV-1；首月以前述100万元初始资金为分母。月度收益复利乘积等于全期收益；不能直接相加。2026-09仅截至21日。\par}

{\small
\begin{longtable}{@{}lrrrr@{}}
\caption{2026年风险仓位与双袖标准成本月度收益（\%）}\label{tab:months-risk}\\
\toprule
月份 & E10 & R01-D & R01-S & 双袖组合 \\
\midrule\endfirsthead
\multicolumn{5}{c}{表\thetable（续）}\\\toprule
月份 & E10 & R01-D & R01-S & 双袖组合 \\
\midrule\endhead
2026-01 & -1.46\% & -0.38\% & -0.30\% & 1.91\% \\
2026-02 & 0.00\% & 0.00\% & 0.17\% & 2.00\% \\
2026-03 & -1.14\% & -2.33\% & -0.94\% & -4.31\% \\
2026-04 & 3.28\% & 0.65\% & 0.91\% & 1.97\% \\
2026-05 & -0.83\% & -2.62\% & -1.29\% & -2.41\% \\
2026-06 & 0.12\% & -0.44\% & -0.15\% & -0.98\% \\
2026-07 & -6.22\% & 0.33\% & -1.77\% & -2.05\% \\
2026-08 & 0.90\% & 1.19\% & 0.92\% & 0.84\% \\
2026-09 & -0.01\% & 0.08\% & -0.08\% & -1.15\% \\
\bottomrule\end{longtable}}
{\footnotesize\noindent 注：区间与表\ref{tab:months}相同。两类R01使用共同股票篮子；固定标量来自开发期校准，不依据2026收益重校准。双袖按真实子账户NAV求和。\par}
各候选采用原固定参数。E00预测与原账户复用，E10/P20各补9次拟合；七组候选共16个账户实例，双袖各情景含两子账户。逐笔费用、滑点、现金/NAV及训练成熟切点检查通过。

RG01的2026账户未完成；SQ01本轮已补固定延伸，结果另节列示；NLP与财务2026完整账户也未完成。缺失结果不以0或推算值替代。2026旧候选比较是在看过E00结果后进行的固定事后对照，不是完全盲测。


\section{新闻主题分析的评价目标与结果}
新闻任务最初混用“标题含证券代码及事件关键词”的个股筛法，随后按研究目标改为六类宏观/产业主题：黄金、石油、科技、航天、宏观风险及流动性。模型复用供应商片段、49篇已核日期正文和11篇政策原文，采用固定字符TF-IDF原型与金融上下文。

唯一固定预测对象为沪深300价格指数五交易日开盘收益。2025年1—5月91个成熟日期训练，6月隔离，7—12月120日期测试；发布日期后第二交易日生效，训练标准化，固定Ridge alpha=1，共三次拟合。Ridge正则能降低估计方差，但并不保证信息有预测价值\cite{ref4}。

表\ref{tab:nlp}给出对应路线或模型的比较结果。


{\small
\begin{longtable}{@{}lrr@{}}
\caption{宏观主题模型的指数五日预测误差}\label{tab:nlp}\\
\toprule
模型 & MSE & MAE \\
\midrule\endfirsthead
\multicolumn{3}{c}{表\thetable（续）}\\\toprule
模型 & MSE & MAE \\
\midrule\endhead
价格基线 & 0.001084312 & 0.02796864 \\
价格＋当前主题 & 0.001392287 & 0.03197059 \\
价格＋延迟20日主题 & 0.000973156 & 0.02505925 \\
零收益参考 & 0.000374764 & 0.01621573 \\
\bottomrule\end{longtable}}
误差单位分别是收益小数的平方与收益小数。当前主题模型差于价格基线；延迟版略好，但三者都差于零收益预测。91日期估计45个变量较不稳定，主题覆盖指示未被完全隔离，因此没有语义的独立贡献证明。预测误差也不等于已完成文本驱动交易账户。

新闻历史首发、版本不可变与完整覆盖未证明，检索片段有选择与截断偏差。新闻可能已被价格吸收，也可能与目标期限不匹配；这些属于解释假设。本实验没有直接预测原油、黄金或行业相对收益，没有证明持续热点交易无效。2026新闻账户未做。


\section{失败原因：事实、解释与尚未完成}
事实层：E00在2025第三季度排序IC约-0.0897，持仓beta与波动分位偏低；同期同池高beta、高波动与趋势端的后续五日收益相对更强。该描述支持当季风格错配，不能推成因果或相加归因。第四季度同一解释并不充分。

机制层：短期排序与慢财务、长期叙事期限不一致；模型适应变化的速度不足；共同量价输入造成信息冗余；区制识别与获利状态不一致；风险控制保护资本却不改变排序。这些合理机制需要专门检验，不能全部列为确定根因。

信息层：公告版本链与宏观历史发布值未齐全，新闻片段未覆盖全部重要事件。模型没有这些输入就不会自动分析地缘冲突、产业订单或油价供需；加入一个HMM或序列算法不会凭空补齐缺失的经济信息。

目标层：前期股票实验以沪深300长仓与现金为主；后期扩展ETF风险来源。公开多资产、多空策略的约束仍与本账户不同。AQR强调不同风格与资产的分散作用\cite{ref1}，并不意味着把几个同源量价模块叠进一个长仓股票池就能复制机构的组合效果。

统计层：截面相关、五日/二十日重叠标签和多次观察减小有效独立样本；更多模型增加事后挑选机会。没有未经查看的足够新样本，也未进行适当多重试验校正，因此本文不声称显著性、因果或普适失效。

工程层：dtype、中文模型路径、前序交易日及旧汇总年度范围曾引起失败，均留档修复。P20标准账户汇总错误只修复汇总、不重跑标准交易或额外拟合。工程失败与经济假设失败分开，不把修复代码当提高策略收益。


以下按实验阶段比较多资产、深度学习与最新配置；各阶段资产池和日期不同，不将其夏普直接横向排名。

\section{多资产配置与阶段选择}
\subsection{先前股票阶段单选的独立结果}
在2025-01-02至2026-09-21的418日连续股票账户中，E00、T01、S01标准夏普分别为0.361、0.555、0.733，历史60日选择基线为0.742；20日主切换规则为0.531，40日敏感性为0.752。20日主规则虽然提高E00表现，仍未胜另外三个简单对照。40日结果高于S01，但仅敏感性且未跑压力，不能事后替换主规则。20日主规则相对S01的夏普差约-0.202，20日移动块区间[-1.108,0.599]跨零。初始持仓与窗口不同，这一连续账户不覆盖旧2026现金重启结果。

\subsection{原四ETF和股票ETF配置}

2025-01-02至2026-09-21，共418日，人民币100万元连续账户；2026为部分年度，已观察探索评价。沪深300ETF代表指数敞口，主动选股同池具有重复市场风险。

表\ref{tab:v3supp1}列示已验收账户结果。
{\footnotesize
\begin{longtable}{@{}lrrrrr@{}}
\caption{既有多资产与股票ETF方案成本后比较}\label{tab:v3supp1}\\
\toprule
方案 & 成本 & 累计收益 & 年化收益 & 夏普 & 最大回撤 \\
\midrule\endfirsthead
\multicolumn{6}{c}{表\thetable（续）}\\\toprule
方案 & 成本 & 累计收益 & 年化收益 & 夏普 & 最大回撤 \\\midrule\endhead
四ETF趋势过滤 & 标准 & 25.19\% & 14.50\% & 1.236 & -11.00\% \\
四ETF固定配置 & 标准 & 25.57\% & 14.72\% & 1.251 & -10.98\% \\
四ETF趋势过滤 & 压力 & 24.40\% & 14.07\% & 1.178 & -11.70\% \\
四ETF固定配置 & 压力 & 24.77\% & 14.27\% & 1.189 & -11.71\% \\
股票ETF阶段选择 & 标准 & 19.60\% & 11.40\% & 0.821 & -12.78\% \\
股票ETF固定组合 & 标准 & 23.33\% & 13.48\% & 1.154 & -8.38\% \\
股票ETF阶段选择 & 压力 & 10.84\% & 6.40\% & 0.502 & -13.65\% \\
股票ETF固定组合 & 压力 & 14.48\% & 8.50\% & 0.826 & -8.87\% \\
\bottomrule\end{longtable}}
原四ETF252日过滤没有胜固定配置；股票ETF40日单选也未胜同候选固定组合。固定组合内部仍选股/轮动，因此结论针对额外选择层，不是所有交易无效。单选账户交易费用反而较少，不能把差异简单归咎费用。

\section{时间注意力：同序列MLP对照}
2020—2022训练、2023验证、2024—2025探索账户485日；20项特征、20交易日窗口、1层32维2头。无新信息源，两模型参数数量不同。

表\ref{tab:v3supp2}列示已验收账户结果。
{\footnotesize
\begin{longtable}{@{}lrrrrrr@{}}
\caption{时间注意力与同序列MLP的种子配对结果}\label{tab:v3supp2}\\
\toprule
模型/种子 & 成本 & 累计收益 & 夏普 & 回撤 & 2024 & 2025 \\
\midrule\endfirsthead
\multicolumn{7}{c}{表\thetable（续）}\\\toprule
模型/种子 & 成本 & 累计收益 & 夏普 & 回撤 & 2024 & 2025 \\\midrule\endhead
ATTN（种子11） & 标准 & 25.21\% & 0.859 & -11.77\% & 12.12\% & 11.67\% \\
ATTN（种子23） & 标准 & 45.84\% & 1.460 & -11.94\% & 26.89\% & 14.94\% \\
ATTN（种子37） & 标准 & 27.15\% & 0.968 & -9.52\% & 19.28\% & 6.60\% \\
MLP（种子11） & 标准 & 20.31\% & 0.703 & -15.63\% & 12.00\% & 7.42\% \\
MLP（种子23） & 标准 & 16.78\% & 0.655 & -12.98\% & 9.38\% & 6.77\% \\
MLP（种子37） & 标准 & 23.87\% & 0.862 & -11.85\% & 11.47\% & 11.13\% \\
ATTN（种子11） & 压力 & 19.08\% & 0.685 & -12.61\% & 9.44\% & 8.81\% \\
ATTN（种子23） & 压力 & 39.48\% & 1.297 & -12.63\% & 24.34\% & 12.18\% \\
ATTN（种子37） & 压力 & 22.37\% & 0.824 & -10.13\% & 16.98\% & 4.61\% \\
MLP（种子11） & 压力 & 14.52\% & 0.537 & -16.17\% & 9.18\% & 4.89\% \\
MLP（种子23） & 压力 & 11.03\% & 0.466 & -14.16\% & 6.44\% & 4.32\% \\
MLP（种子37） & 压力 & 17.55\% & 0.670 & -12.86\% & 8.04\% & 8.80\% \\
\bottomrule\end{longtable}}
注意力三seed标准与压力均胜对应MLP，标准平均夏普差约0.356，压力约0.378。标准配对块区间约[-0.015,0.829]跨零，种子23贡献较大；三seed不是独立市场样本，均值不是可交易组合。同期滚动E00夏普1.595仍较高，但静态/滚动更新不同，不能据此隔离架构影响。队友已返回静态匹配树、排序融合和跨资产注意力结果；补证范围与限制另节列示。

\section{数据、未完成与提交边界}
A新增三ETF：513100纳指100、510500中证500、159980工业金属期货。本轮中证500已进入四ETF账户；纳指和有色期货因尾段权益证据缺口未准入，不能列作已完成的轮动策略。513100早期份额事件与本次2024Q4预热/2025起账户准入须分别说明。市场融资余额不是个股资金，队员标题沪深与含北交所实际口径不一致；两日发布延迟尚是假设。北向缺失不填零、龙虎榜事后成功率不作训练特征。

现金始终允许，零计息为基线口径；债券配置不是熊市强制买入。NLP现有沪深300五日预测负结果不证明行业/油价新闻普遍无效。所有新结果仍是已观察探索，不宣称实盘或独立盲测。

\section{新四ETF配置与价格现金门}
2025-01-02至2026-09-21，418日，100万元。四模块分别为中证500ETF、标普500ETF、黄金ETF、国债ETF，代码依次为510500、513500、518800、511010，目标资本各25\%。C01固定配置；S02在40日节点仅当60日累计总收益除以波动的机会分数为正才保留该模块，否则现金。现金零计息；不能将该60日价格机会规则描述为120日账户夏普选择。

表\ref{tab:v3supp3}列示已验收账户结果。
{\footnotesize
\begin{longtable}{@{}lrrrrr@{}}
\caption{新四ETF固定配置与价格现金门}\label{tab:v3supp3}\\
\toprule
方案 & 成本 & 累计收益 & 年化收益 & 夏普 & 最大回撤 \\
\midrule\endfirsthead
\multicolumn{6}{c}{表\thetable（续）}\\\toprule
方案 & 成本 & 累计收益 & 年化收益 & 夏普 & 最大回撤 \\\midrule\endhead
新四ETF固定配置 & 标准 & 30.05\% & 17.16\% & 1.276 & -12.84\% \\
新四ETF固定配置 & 压力 & 29.70\% & 16.97\% & 1.259 & -12.86\% \\
新四ETF价格现金门 & 标准 & 23.69\% & 13.68\% & 1.256 & -10.98\% \\
新四ETF价格现金门 & 压力 & 23.09\% & 13.34\% & 1.229 & -10.98\% \\
\bottomrule\end{longtable}}
S02标准夏普差-0.0207、压力-0.0301，均未过预设增量门；回撤略小不构成以夏普为目标的升级。40日配对标准区间[-0.573,0.550]跨零。固定配置保留为多资产候选，不能以较高累计收益证明未来必胜。纳指、有色期货因评价尾段权益证据尚未闭合在结果前排除，本轮不检验模块内轮动；股票S01适配与F01融资覆盖受阻也未执行。

表\ref{tab:v3supp4}列示已验收账户结果。
{\footnotesize
\begin{longtable}{@{}lrrr@{}}
\caption{新四ETF方案分年收益}\label{tab:v3supp4}\\
\toprule
方案 & 成本 & 2025收益 & 2026截至9月21日收益 \\
\midrule\endfirsthead
\multicolumn{4}{c}{表\thetable（续）}\\\toprule
方案 & 成本 & 2025收益 & 2026截至9月21日收益 \\\midrule\endhead
新四ETF固定配置 & 标准 & 25.11\% & 3.95\% \\
新四ETF固定配置 & 压力 & 24.93\% & 3.82\% \\
新四ETF价格现金门 & 标准 & 19.38\% & 3.61\% \\
新四ETF价格现金门 & 压力 & 19.19\% & 3.27\% \\
\bottomrule\end{longtable}}
年际净值连续，不重置现金及持仓；2026是部分年度。逐资产损益恒等式、现金/份额/权益/NAV重建和成交检查通过；历史日交易状态非完整档案，国债九月分红仍使用次级存证，原公告转载与原PDF分别标明。

\section{序列SQ01的2026固定延伸}
2026-01-05至9月21日，共175日；复用冻结32维ESN状态，36个月训练窗及成熟五日标签，1—9月各一次Ridge读出拟合。51893个预测键与E00一致，128键冷启动回退，参数不搜索。

表\ref{tab:v3supp5}列示已验收账户结果。
{\footnotesize
\begin{longtable}{@{}lrrrrr@{}}
\caption{SQ01序列路线2026延伸结果}\label{tab:v3supp5}\\
\toprule
方案 & 成本 & 累计收益 & 夏普 & 最大回撤 & 费用（元） \\
\midrule\endfirsthead
\multicolumn{6}{c}{表\thetable（续）}\\\toprule
方案 & 成本 & 累计收益 & 夏普 & 最大回撤 & 费用（元） \\\midrule\endhead
SQ01 & 标准 & -6.15\% & -0.753 & -11.08\% & 5297.62 \\
SQ01 & 压力 & -6.86\% & -0.849 & -11.47\% & 10543.98 \\
\bottomrule\end{longtable}}
对同窗E00的标准/压力收益差为-1.85/-0.86个百分点，夏普差-0.301/-0.185；费用更低但表现更弱，不支持替换。日期、键、费用、现金重放与NAV恒等式通过；独立检查未从原事件逐笔重建持仓批次，验收范围不夸大。RG01的2026账户仍未做，不能与SQ01一起写成已经完成。


\section{最终候选：相对选股与跨资产风险预算}
投资策略分为两个层次。E00在历史沪深300成分内选择相对看好的股票；资产层再决定股票、标普500、黄金和国债四个模块承担多少风险。境内标普500ETF引入美国大盘股票风险，黄金引入不同的商品与避险驱动，国债引入利率风险。四者并非独立，相关性会变化，因此分散效果须由实际账户检验。

固定配置对照每20交易日恢复四模块各25\%的目标权重，股票模块仍每5交易日按E00排名检查持仓。它包含交易与再平衡，不能解释为静态买入持有。动态候选RISK-B1在同一20日节点使用截至当日收盘的60日收益协方差，次日开盘执行；协方差向其对角阵收缩50\%，降低短窗口相关性估计的不稳定。

股票、标普500、黄金、国债的风险贡献目标依次为45\%、30\%、20\%、5\%，这是风险份额而非资本权重。风险代理的股票列为中证500，实际股票账户执行E00，代理错配保留为局限。令$\widehat\Sigma_t$为收缩后的协方差，定义
\begin{equation}
\sigma_p(w)=\sqrt{w^\top\widehat\Sigma_t w},\qquad
RC_i(w)=\frac{w_i(\widehat\Sigma_t w)_i}{\sigma_p(w)}.
\end{equation}
每20日求解
\begin{equation}
\min_w\sum_{i=1}^4\left(\frac{RC_i(w)}{\sum_jRC_j(w)}-b_i\right)^2
+0.02\left(\sum_i|w_i-w_{i,t-1}|\right)^2,
\end{equation}
满足$w_i\ge0$、$\sum_iw_i=1$，四模块权重上限为0.8、0.7、0.6、0.6。数值求解采用SLSQP，失败时保持上期目标；这不是每期精确实现风险份额的保证。现金始终允许，由整手、成本、未成交与目标余量形成；本候选没有额外训练独立的现金机会预测器。

另一个风险预算B2及固定公式趋势倾斜作为候选对照。B1在2023年开发验证中的零现金夏普2.837高于B2的2.776，因此在两风险候选之间选用B1。2023年多次用于研发，仅一个已查看的开发年度，2024—2025又曾反复观察，故最终结果属于探索性比较，不能称全新盲测或统计显著胜出。
\subsection{固定比例搜索、收益预测与实际配置}
固定25\%并未被假定为最优。曾在2020—2022年代理路径上评价132组预设、去重的可行固定权重，以各年度夏普均值减0.5倍年度夏普标准差评分，并排除任一年非正收益或指标无效的候选。选出两组差异较大的方案后，只在2023年真实共享账户中验证：两者夏普2.292和2.267，均低于固定25\%的2.662，因此该轮保留简单对照。该有限搜索不证明不存在更好的比例；训练代理采用近似外部换手成本，与真实整手账户不同。

还比较了20日收益预测驱动的Ridge和小型MLP配置，训练2020—2022年、开发验证2023年。在2024—2025共享账户中，标准夏普分别约1.788和1.811，低于固定配置2.642。该轮标签使用原价开盘收益，股票预测目标采用指数代理；它没有完全覆盖分红及E00自身的条件收益。预测误差、代理错配及仓位转换共同作用，不能把不胜对照归结为某一种网络无用。

表\ref{tab:finalmain}给出最终两个主方案及趋势倾斜对照；表\ref{tab:finalvalidation}列出开发验证，表\ref{tab:finalfees}列出逐笔记录合计。所有主结果来自共享现金、股票整手及ETF交易的账户，而不是把四条净值直接求平均。
{\small
\begin{longtable}{@{}lrrrrr@{}}
\caption{股票与三ETF共享现金账户比较（2024—2025年）}\label{tab:finalmain}\\
\toprule
方案 & 成本 & 累计收益 & 年化收益 & 夏普 & 最大回撤 \\
\midrule\endfirsthead
\multicolumn{6}{c}{表\thetable（续）}\\\toprule
方案 & 成本 & 累计收益 & 年化收益 & 夏普 & 最大回撤 \\
\midrule\endhead
固定25\% & 标准 & 46.13\% & 21.78\% & 2.642 & -5.55\% \\
RISK-B1 & 标准 & 27.80\% & 13.59\% & 2.858 & -3.38\% \\
趋势倾斜 & 标准 & 40.95\% & 19.52\% & 2.393 & -4.96\% \\
固定25\% & 压力 & 44.40\% & 21.03\% & 2.556 & -5.60\% \\
RISK-B1 & 压力 & 27.03\% & 13.24\% & 2.742 & -3.44\% \\
趋势倾斜 & 压力 & 38.03\% & 18.23\% & 2.254 & -5.05\% \\
Ridge配置 & 标准 & 39.55\% & 18.91\% & 1.788 & -8.33\% \\
MLP配置 & 标准 & 39.64\% & 18.94\% & 1.811 & -9.42\% \\
Ridge配置 & 压力 & 36.12\% & 17.38\% & 1.686 & -8.47\% \\
MLP配置 & 压力 & 38.19\% & 18.30\% & 1.775 & -9.41\% \\
\bottomrule\end{longtable}}
{\footnotesize\noindent 注：2024-01-02至2025-12-31，485日，起始100万元；股票模块执行E00，ETF为513500、518800、511010。标准及压力各自重放账户，回撤包含初始资金；零现金夏普采用日收益样本标准差、252日约定。压力费用及滑点加倍；国债仍承担价格风险。固定25\%指目标权重，不是每日实际权重，也不是买入后永久不调仓。\par}
{\small
\begin{longtable}{@{}lrrrrr@{}}
\caption{2023年开发验证账户（242日、标准成本）}\label{tab:finalvalidation}\\
\toprule
方案 & 成本 & 累计收益 & 年化收益 & 夏普 & 最大回撤 \\
\midrule\endfirsthead
\multicolumn{6}{c}{表\thetable（续）}\\\toprule
方案 & 成本 & 累计收益 & 年化收益 & 夏普 & 最大回撤 \\
\midrule\endhead
固定25\% & 标准 & 12.52\% & 13.07\% & 2.662 & -2.50\% \\
RISK-B1 & 标准 & 8.49\% & 8.85\% & 2.837 & -1.74\% \\
RISK-B2 & 标准 & 7.81\% & 8.15\% & 2.776 & -1.59\% \\
趋势倾斜 & 标准 & 12.41\% & 12.96\% & 2.614 & -2.94\% \\
\bottomrule\end{longtable}}
{\small
\begin{longtable}{@{}lrrrrr@{}}
\caption{2024—2025年共享账户成交及成本}\label{tab:finalfees}\\
\toprule
方案 & 成本 & 成交笔数 & 成交额（万元） & 费用（元） & 股息税（元） \\
\midrule\endfirsthead
\multicolumn{6}{c}{表\thetable（续）}\\\toprule
方案 & 成本 & 成交笔数 & 成交额（万元） & 费用（元） & 股息税（元） \\
\midrule\endhead
固定25\% & 标准 & 2422 & 749.65 & 13,906.71 & 1,729.18 \\
固定25\% & 压力 & 2412 & 738.64 & 27,690.25 & 1,699.43 \\
RISK-B1 & 标准 & 2082 & 530.38 & 11,667.54 & 788.49 \\
RISK-B1 & 压力 & 2023 & 508.27 & 22,670.30 & 798.97 \\
\bottomrule\end{longtable}}
{\small
\begin{longtable}{@{}lrrrr@{}}
\caption{最终主方案年度收益（连续账户，\%）}\label{tab:finalannual}\\
\toprule
年度 & 固定25\%标准 & B1标准 & 固定25\%压力 & B1压力 \\
\midrule\endfirsthead
\multicolumn{5}{c}{表\thetable（续）}\\\toprule
年度 & 固定25\%标准 & B1标准 & 固定25\%压力 & B1压力 \\
\midrule\endhead
2024 & 24.64\% & 17.55\% & 23.75\% & 16.98\% \\
2025 & 17.23\% & 8.72\% & 16.68\% & 8.59\% \\
\bottomrule\end{longtable}}
{\small
\begin{longtable}{@{}lrrrr@{}}
\caption{最终主方案月度收益（连续账户，\%）}\label{tab:finalmonthly}\\
\toprule
月份 & 固定25\%标准 & B1标准 & 固定25\%压力 & B1压力 \\
\midrule\endfirsthead
\multicolumn{5}{c}{表\thetable（续）}\\\toprule
月份 & 固定25\%标准 & B1标准 & 固定25\%压力 & B1压力 \\
\midrule\endhead
2024-01 & 1.33\% & 1.22\% & 1.27\% & 1.15\% \\
2024-02 & 2.58\% & 1.95\% & 2.47\% & 1.81\% \\
2024-03 & 3.80\% & 3.51\% & 3.69\% & 3.39\% \\
2024-04 & 0.17\% & 0.08\% & 0.04\% & -0.08\% \\
2024-05 & 1.41\% & 0.86\% & 1.40\% & 1.04\% \\
2024-06 & 2.14\% & 1.67\% & 2.12\% & 1.73\% \\
2024-07 & 0.94\% & 0.59\% & 0.88\% & 0.56\% \\
2024-08 & 0.00\% & -0.03\% & -0.02\% & -0.08\% \\
2024-09 & 3.13\% & 2.09\% & 3.11\% & 2.07\% \\
2024-10 & 3.69\% & 2.72\% & 3.61\% & 2.64\% \\
2024-11 & 2.69\% & 1.14\% & 2.65\% & 1.12\% \\
2024-12 & 0.45\% & 0.55\% & 0.38\% & 0.50\% \\
2025-01 & 1.15\% & 0.83\% & 1.08\% & 0.80\% \\
2025-02 & -1.14\% & -1.09\% & -1.18\% & -1.11\% \\
2025-03 & 0.77\% & 0.36\% & 0.74\% & 0.31\% \\
2025-04 & 1.68\% & 1.37\% & 1.60\% & 1.32\% \\
2025-05 & 1.55\% & 0.67\% & 1.47\% & 0.82\% \\
2025-06 & 1.45\% & 0.83\% & 1.46\% & 1.01\% \\
2025-07 & 1.85\% & 0.93\% & 1.79\% & 0.90\% \\
2025-08 & 0.71\% & 0.25\% & 0.68\% & 0.21\% \\
2025-09 & 3.55\% & 1.99\% & 3.50\% & 1.92\% \\
2025-10 & 3.82\% & 2.07\% & 3.80\% & 1.99\% \\
2025-11 & 0.75\% & 0.24\% & 0.75\% & 0.22\% \\
2025-12 & -0.04\% & -0.00\% & -0.07\% & -0.07\% \\
\bottomrule\end{longtable}}
{\footnotesize\noindent 注：表\ref{tab:finalannual}与表\ref{tab:finalmonthly}按上年/上月最后净值衔接，首期以100万元为分母；跨年不重置持仓，月度收益复利相乘得到全期收益。表中实际收益均为各情景账户净收益。\par}
RISK-B1的标准夏普相对固定配置增加约0.216，压力增加约0.186；最大回撤分别由5.55\%降至3.38\%、5.60\%降至3.44\%。代价是标准累计收益由46.13\%降至27.80\%，年化由21.78\%降至13.59\%。这符合降低组合波动和改变资产暴露的经济机制，但不证明新增股票选股alpha。按本项目后期确定的夏普优先目标，B1作为最终候选，固定配置保留为强对照；不同目标的投资者可能作不同选择。

RISK-B1在2024和2025年的标准收益分别为17.55\%、8.72\%，2025没有出现股票基线那样明显的收益落差，但这一事实不足以证明宏观择时成功。平均目标资本权重约为股票14.16\%、标普14.87\%、黄金14.54\%、国债56.42\%；对应平均实际股票仅8.08\%，ETF合计81.14\%，现金10.78\%。低股票预算在100万元、80目标股票及整手约束下留下未部署资金，必须把这种执行影响纳入策略定位。低波动国债取得较大资本权重，是风险份额目标的数学结果，并非“判断熊市后强制买债”。

固定配置的平均实际股票17.99\%、ETF合计71.89\%、现金10.13\%，同样偏离目标资本权重。压力情景改变费用与可成交数量，持仓路径随之变化；不能用标准路径简单扣除两倍费用代替压力回放。ETF卖出记录曾遗漏于导出的成交统计，修复后重新回放，净值不变、成交数及费用统计更正；本报告采用更正后的数值。

最终候选已补充2026年175日现金重启账户，结果另表列示；尚无继承2025持仓的2024—2026连续账户。股票候选、旧四ETF过滤及深度学习代理结果与最终共享账户分别披露。回测中ETF使用境内交易价格，QDII溢折价、汇率、交易时区及流动性风险会影响可执行收益；日线滑点假设不能精确描述盘口。

\subsection{深度学习返回包的证据分层}
静态训练2020—2022、开发验证2023的匹配LightGBM（N01）标准夏普约0.847；与注意力的比较较滚动E00更接近相同更新方式。E00与注意力百分位50/50融合（N02）三种子标准夏普约0.884、1.327、0.918，仅为摘要支持，未完整独立重放。

四ETF跨资产试验（N03）的资产对象为沪深300、标普500、黄金、国债ETF，2022训练仅161日、2023验证222日，与本节E00主动股票对象不同。2024—2025标准夏普：固定配置约2.439，MLP三种子约2.913、2.327、2.321，注意力约1.219、1.461、1.826。单一种子成功不能证明稳定模型优势。十四个标准/压力导出账户的收益、回撤及费用合计已经独立算术核验；补存权重是再训练得到，原预测和完整标准化工件未全部交付，不能声称从原输入完整复现了网络决策。该轮原价标签被称为总收益、年度衔接及波动字段命名存在口径问题，故不作为正式候选的核心证据。

E00股票与三ETF的N04摘要曾缺账户明细；后续本机按共享现金引擎重新回放，形成表\ref{tab:finalmain}固定配置及B1的独立账本。正式主结果以本机更正后账本为准，不直接引用返回包最优摘要。最终主方法无需PyTorch；深度学习保留为有限样本的比较研究。

时间注意力的2026外部返回账户虽已收到，但预测缺失时排序回退占槽与冻结规则不一致，未纳入正式架构胜负判断。收到输出文件不等于通过执行验收。
\subsection{2026年最终候选的冻结规则补充检验}
为检验最终配置而不是以旧股票或四ETF实验代替，本次按既定规则补做RISK-B1与同对象固定配置的标准、压力账户。区间为2026年1月5日至9月21日，共175个交易日，各账户以100万元现金重新起步，不承接2025持仓；因此它不是2024—2026连续净值，也不是2026全年结果。E00使用已完成的九次月度模型与冻结评分，风险预算参数、60日风险窗口及20日外层再平衡规则不因这次结果改动。后月训练可使用此前已经成熟的2026标签，未使用未来未成熟标签。
{\small
\begin{longtable}{@{}lrrrr@{}}
\caption{2026年最终共享账户检验（175日、现金起点）}\label{tab:final2026}\\
\toprule
方案 & 成本 & 区间收益 & 夏普 & 最大回撤 \\
\midrule\endfirsthead
\multicolumn{5}{c}{表\thetable（续）}\\\toprule
方案 & 成本 & 区间收益 & 夏普 & 最大回撤 \\
\midrule\endhead
RISK-B1 & 标准 & 2.96\% & 0.446 & -8.76\% \\
固定25\% & 标准 & 0.52\% & 0.122 & -10.22\% \\
RISK-B1 & 压力 & 2.78\% & 0.433 & -8.72\% \\
固定25\% & 压力 & 0.26\% & 0.091 & -10.18\% \\
\bottomrule\end{longtable}}
{\small
\begin{longtable}{@{}lrrrrr@{}}
\caption{2026年账户执行与平均实际暴露}\label{tab:final2026cost}\\
\toprule
方案 & 成本 & 成交笔数 & 费用（元） & 现金占比 & 股票占比 \\
\midrule\endfirsthead
\multicolumn{6}{c}{表\thetable（续）}\\\toprule
方案 & 成本 & 成交笔数 & 费用（元） & 现金占比 & 股票占比 \\
\midrule\endhead
RISK-B1 & 标准 & 1064 & 6,281.47 & 15.82\% & 15.78\% \\
固定25\% & 标准 & 963 & 5,571.59 & 15.32\% & 13.62\% \\
RISK-B1 & 压力 & 1060 & 12,407.74 & 39.91\% & 15.85\% \\
固定25\% & 压力 & 957 & 10,955.35 & 39.27\% & 13.65\% \\
\bottomrule\end{longtable}}
{\small
\begin{longtable}{@{}lrrrr@{}}
\caption{2026年最终共享账户月度收益}\label{tab:final2026monthly}\\
\toprule
月份 & B1标准 & 固定标准 & B1压力 & 固定压力 \\
\midrule\endfirsthead
\multicolumn{5}{c}{表\thetable（续）}\\\toprule
月份 & B1标准 & 固定标准 & B1压力 & 固定压力 \\
\midrule\endhead
2026-01 & 2.35\% & 3.47\% & 2.59\% & 3.69\% \\
2026-02 & -1.19\% & -0.78\% & -0.88\% & -0.50\% \\
2026-03 & -4.58\% & -5.12\% & -4.79\% & -5.30\% \\
2026-04 & 4.24\% & 2.84\% & 4.08\% & 2.70\% \\
2026-05 & 1.06\% & 0.35\% & 1.02\% & 0.26\% \\
2026-06 & -2.86\% & -4.01\% & -2.88\% & -4.07\% \\
2026-07 & 0.79\% & 0.83\% & 0.66\% & 0.71\% \\
2026-08 & 4.27\% & 4.19\% & 4.17\% & 4.10\% \\
2026-09 & -0.80\% & -0.84\% & -0.86\% & -0.90\% \\
\bottomrule\end{longtable}}
{\footnotesize\noindent 注：月度收益按前月最后净值衔接；一月从初始资金开始，九月截至21日。费用为逐笔显式交易费用，滑点已反映于成交价格，股息税另计。夏普仍采用零现金基准及252日约定，不将部分年度收益冒充全年收益。\par}

标准情景下，B1比固定配置多取得约2.44个百分点区间收益，夏普高约0.324，回撤浅约1.46个百分点；压力情景结论方向一致。但B1的标准夏普仅0.446，最大回撤扩大至8.76\%，显著弱于2024—2025的2.858和3.38\%。因此，跨资产分散与风险预算具有相对改善价值，尚不能保证所有市场阶段都获得高夏普。2026已经查看过背景与部分分支，仍属于事后补充检验，不能宣称全新盲测。

2026的平均目标资本权重约为股票27.46\%、标普31.99\%、黄金15.46\%、国债25.09\%，不同于2024—2025国债约56.42\%的均值；B1并非固定重仓国债。标准实际股票约15.78\%、ETF合计68.39\%、现金15.82\%。压力成本改变整手成交与持仓路径，实际现金升至39.91\%，所以压力回撤略浅不代表成本增加有利。现金未单独计息。

账户内部净资产恒等式、费用合计及非负现金检查通过。代码包提供冻结预测重放入口与参考账本；复跑可核对上述账户，但不把这种执行验收等同所有扩展模型重新训练或金融假设完全正确。尚未补做继承2025持仓的连续账户、独立2026新闻交易及完整日内成交模拟。国债9月权益事件的初级公告来源仍有缺口，来源限制随输入清单披露。

\section{结论、可复现性与证据索引}
结论一：量价股票基线具有显著年度分化。财务、趋势、隐状态、序列及新闻路线提供局部互补或负结果，未达到各轮设定的升级目标；这只限制本次实现，不否定方法的一般用途。

结论二：扩展到不同经济风险来源后，固定配置形成强对照。风险预算B1在2024—2025年提高成本后夏普、降低回撤，同时减少绝对收益，符合后期夏普优先目标。收益改善来自风险承担、分散与执行的共同变化，不能包装成新增股票alpha或已经实现的新闻宏观择时。

结论三：最终候选仍需新的未查看时期、不同资金规模和风险代理检验。2026最终组合的175日现金重启结果已列明，B1仍优于固定配置，但绝对夏普下降、回撤扩大；股票与旧ETF分支分别保留，连续跨年检验仍未完成。2023选择和反复查看2024—2025均带有选择偏差，不宣称全局最优、因果或统计显著。

\subsection*{交付与复现说明}
正式提交为报告PDF、代码ZIP、数据ZIP三个文件，每个独立不超过100MB；代码与数据共同解压后，按代码说明重放主候选及固定配置的2023开发验证、2024—2025及2026部分年度标准/压力共享账户。E00训练提供独立入口，标签成熟、数据版本和训练范围以数据清单及代码说明为准。环境另行安装；冻结预测回放、模型重训与旧分支结果重汇总是三种不同的复现范围，不能混称全部实验重训通过。

本文为评价结果，完整成交与输入SHA在代码、数据的清单中对应。主结果采用更正ETF卖出记录后的账本；扩展研究保留方法、参数、汇总及来源索引，不把未封装的完整研究档案伪装成轻量提交可复现。新闻及供应商资料按比赛提交范围使用，不包含访问凭证。
\section*{指标定义}
对日净值$V_t$及起始净值$V_0$，累计收益、年化收益及最大回撤定义为
\begin{align}
R&=V_N/V_0-1, & R_{\mathrm{ann}}&=(V_N/V_0)^{252/N}-1,\\
\mathrm{MDD}&=\min_{0\le t\le N}\left\{\frac{V_t}{\max_{0\le s\le t}V_s}-1\right\}.
\end{align}
日收益为$r_t=V_t/V_{t-1}-1$。以零现金收益为参考，年化波动与Sharpe为
\begin{equation}\sigma_{\mathrm{ann}}=s(r)\sqrt{252},\qquad S_0=\sqrt{252}\,\overline r/s(r),\end{equation}
其中$s(r)$为样本标准差（自由度$N-1$）。短期样本的年化指标仅用于描述；2026累计收益不改写为全年实际收益。
\begingroup\raggedright\begin{thebibliography}{9}
\bibitem{ref1} Israel R, Maloney T. Understanding Style Premia[EB/OL]. AQR, 2014. \url{https://www.aqr.com/Insights/Research/Journal-Article/Understanding-Style-Premia}.
\bibitem{ref2} Microsoft. Qlib: An AI-oriented Quantitative Investment Platform[EB/OL]. \url{https://github.com/microsoft/qlib}.
\bibitem{ref3} hmmlearn Developers. Tutorial[EB/OL]. \url{https://hmmlearn.readthedocs.io/en/stable/tutorial.html}.
\bibitem{ref4} scikit-learn Developers. Ridge[EB/OL]. \url{https://scikit-learn.org/stable/modules/generated/sklearn.linear_model.Ridge.html}.
\end{thebibliography}\endgroup
\end{document}
